% ===================================================================
%  TAULUKKO N:o 10 — Meri. — Satama.
%  Traduction 2026 d'apres texto/io, controlee sur texto/fr et
%  texto/sv. Consigne : voir texto/fi/10-taulukko-01.tex.
%
%  Le francais decale sa numerotation d'un cran a partir de l'alinea
%  8 ; c'est une coquille de l'atelier francais, et la colonne suit
%  l'ido, qui compte juste.
% ===================================================================

%%K t10-apar-1 apar
\VUsaut{4.0mm}
\VUtitre{15.0pt}{60}{TAULUKKO N:o 10}
\VUsaut{3.0mm}
\VUpk{11.4pt}{40}{Meri. --- Satama.}
\VUsaut{3.0mm}

%%K t10-01-1 p
1. --- Kesällä, samalla kun jotkut etsivät rauhaa ja viileyttä vuorelta, toiset
suosivat matkaa rannikolle. Niin teimme viime vuonna, sillä pidämme kovasti
\VUgras{valtamerestä} \textsuperscript{(1)}.

%%K t10-02-1 p
2. --- Mitä minuun tulee, tunnen suurta mielihyvää katsellessani sitä valtavaa
vesialaa, joka ulottuu aina \VUgras{taivaanrantaan} \textsuperscript{(2)} asti,
ja nähdessäni valtavien \VUgras{höyrylaivojen} \textsuperscript{(3)} lähtevän
matkaan pitkälle \VUgras{ylimenolle} matkustajineen, jotka menevät Amerikkaan.

%%K t10-03-1 p
3. --- Siksi isäni, joka tuntee taipumukseni, valitsee aina hyvin rauhallisen
rannan lomien viettoon, joka on jonkin suuren \VUgras{sotasatamamme} luona.

%%K t10-03-2 p
Viimeisen oleskelumme aikana \VUgras{laivasto} tuli ankkuroimaan valtavan
\VUgras{aallonmurtajan} \textsuperscript{(4)} taakse, joka sulkee ja suojaa
\VUgras{sotasataman} \textsuperscript{(5)}, ja onnistuin toiveeni mukaan
ihailemaan eri \VUgras{sota-aluksia}, pientä \VUgras{sukellusvenettä}
\textsuperscript{(10)}, joka sukeltaa ja katoaa veden alle, ja myös valtavaa
\VUgras{panssarilaivaa} \textsuperscript{(9)}, joka tähän asti pelkäsi miltei
vain \VUgras{torpedoveneen} \textsuperscript{(8)} hyökkäystä.

%%K t10-tit-2 sub
\VUsaut{3.0mm}
\VUpk{10.2pt}{40}{Pienet alukset.}
\VUsaut{3.0mm}

%%K t10-04-1 p
4. --- Tämä viimeksi mainittu osasi jo hyvin \VUgras{taitavasti} löytää
paksun metallipanssarin heikon kohdan pannakseen siihen vaarallisen
\VUgras{torpedon}, jonka räjähdys aiheuttaa aluksen tuhon, mutta kuitenkin sitä
oli vähemmän pelättävä kuin sukellusvenettä, sillä hävittäjät ajavat sitä
helposti takaa.

%%K t10-05-1 p
5. --- Saatoin myös nähdä nopeat panssaroidut \VUgras{risteilijät}
\textsuperscript{(6)}, miltei yhtä haavoittumattomat kuin varsinainen
panssarilaiva, useita \VUgras{kuljetusaluksia} \textsuperscript{(12)}, kolme tai
neljä \VUgras{tykkivenettä} \textsuperscript{(7)} ja lopuksi \VUgras{fregatin}
\textsuperscript{(11)}. \VUgras{Tuon laivaston näytelmä, kun se liikehtii
nostaakseen ankkurin, on kaikkein mielenkiintoisin.}

%%K t10-06-1 p
6. --- Minkä eron rakenteessa suurten teräsmöhkäleiden ja tyylikkäiden
\VUgras{kolmimastoisten} tai sulavien \VUgras{purjealusten}
\textsuperscript{(13)} välillä, joita vielä tänäänkin käytetään
kauppalaivastossa ja joiden näin tulevan satamaan, täysin purjein, kun kävelin
\VUgras{laiturilla} \textsuperscript{(14)}. Tosin nämä menettivät paljon
eduistaan, \VUgras{kun} \VUgras{tuuli} oli vastainen. Niiden täytyi laskea alas
kaikki purjeet päästäkseen takaisin altaaseen, ja \VUgras{hinaajaa}
\textsuperscript{(16)} tarvittiin hakemaan ne ja viemään ne, kuin pelkät
\VUgras{proomut} \textsuperscript{(17)}, laiturille.

%%K t10-tit-3 sub
\VUsaut{3.0mm}
\VUpk{10.2pt}{40}{Kauppasatama.}
\VUsaut{3.0mm}

%%K t10-07-1 p
7. --- Se, mikä on mielenkiintoista siinä satamassa, josta puhun, on se, ettei se
sisällä vain sotasatamaa, keinotekoisesti rakennettuna jättiläismäisin töin, vaan
myös ihmeellisen \VUgras{kauppasataman}, joka on suuren joen
\VUgras{suistossa}.

%%K t10-08-1 p
8. --- Maakieleke, joka erottaa kaksi allasta, on linnoitettu ja puolustettu
\VUgras{pattereiden} ja \VUgras{bastionien} rivillä, jotka työntyvät mereen.
Tarvitaan erityinen lupa kävelläkseen \VUgras{valleilla}
\textsuperscript{(15)}. Sain sen helposti setäni kautta, joka on insinööri
\VUgras{telakoilla} \textsuperscript{(18)}, ja menin katsomaan niitä aluksia,
joita rakennetaan laajojen vajojen alla.

%%K t10-09-1 p
9. --- Olin jopa läsnä panssarilaivan vesillelaskussa, ja muistan sen
liikuttavan hetken, jonka koko väkijoukko tunsi, kun valtava möhkäle, jätettynä
oman onnensa nojaan, alkoi liukua kehtomaisella telallaan ja tuli kellumaan joen
vedelle.

%%K t10-10-1 p
10. --- Telakoille mennäkseen kuljetaan \VUgras{harjoituskentän}
\textsuperscript{(19)} yli, jossa varuskunnan sotilaat tulevat joka päivä
harjoittelemaan, kuljetaan rautaisen \VUgras{pienen sillan}
\textsuperscript{(20)} yli, joka yhdistää paraatikentän \VUgras{arsenaaliin}
\textsuperscript{(21)}.

%%K t10-tit-4 sub
\VUsaut{3.0mm}
\VUpk{10.2pt}{40}{Arsenaali ja telakka-altaat.}
\VUsaut{3.0mm}

%%K t10-11-1 p
11. --- Setäni kanssa kävin siinä suuressa rakennuksessa, joka on täynnä
\VUgras{tykkejä} \textsuperscript{(22)}, \VUgras{kuulia}
\textsuperscript{(23)} ja \VUgras{kranaatteja}. Näin myös
\VUgras{metallitehtaan} \textsuperscript{(24)}, jossa valmistetaan
höyrylaivojen \VUgras{höyrykattilat} \textsuperscript{(25)}.

%%K t10-12-1 p
12. --- Aivan vieressä ovat \VUgras{telakka-altaat} \textsuperscript{(26)} ---
korjausaltaat --- ja hyvin merkilliset \VUgras{kelluvat telakat}
\textsuperscript{(27)}, joissa saatoin hyvin läheltä, korjattavana olevassa
aluksessa, nähdä veneen \VUgras{potkurin} \textsuperscript{(28)},
\VUgras{kölin} \textsuperscript{(29)} ja \VUgras{peräsimen}
\textsuperscript{(30)}.

%%K t10-13-1 p
13. --- Kauppasatama on yhtä tutkimisen arvoinen kuin sotasatama, ja vietin monta
tuntia töllistellen laitureilla, joilla \VUgras{siipirataslaivat}
\textsuperscript{(31)} ovat kiinnitettyinä ja jotka kulkevat jokea ylös, ja
katsellen, kuinka \VUgras{mastonosturit} \textsuperscript{(32)}
työskentelevät, jotka, kuin jouset, nostavat valtavia kuormia.

%%K t10-tit-5 sub
\VUsaut{4.0mm}
\VUpk{10.2pt}{40}{Luotsi.}
\VUsaut{3.0mm}

%%K t10-14-1 p
14. --- Ihailin \VUgras{valtamerilaivoja} \textsuperscript{(34)}, jotka lähtivät
redillä, \VUgras{liput} \textsuperscript{(35)} täydessä tuulessa, taitavan
\VUgras{luotsin} \textsuperscript{(36)} ohjaamina, joka ihmeen hyvin osaa
seurata \VUgras{pientä väylää}, joka on merkitty \VUgras{poijuilla}
\textsuperscript{(40)} ja \VUgras{merimerkeillä}, samalla kun hän väistää
\VUgras{proomuja} \textsuperscript{(41)}, joita ohjataan vaivalloisesti niiden
ainoalla nelikulmaisella purjeella. Erotin \VUgras{keulassa}
\textsuperscript{(37)} --- keulapuolella --- toisen luokan \VUgras{hytit}
\textsuperscript{(39)}, ja \VUgras{perässä} \textsuperscript{(38)} ---
peräpuolella --- ensimmäisen luokan matkustajien salongit.

%%K t10-15-1 p
15. --- Joskus olin myös läsnä \VUgras{proomun} \textsuperscript{(42)}
lastauksessa, johon juuri pinottiin tavaroita \VUgras{makasiinista}
\textsuperscript{(43)} ja \VUgras{merenrannan asemalta} \textsuperscript{(44)},
sinne tuotuina lukuisilla junilla \VUgras{laituriradalla}
\textsuperscript{(45)}.

%%K t10-tit-6 sub
\VUsaut{3.0mm}
\VUpk{10.2pt}{40}{Huvisoutu.}
\VUsaut{3.0mm}

%%K t10-16-1 p
16. --- Usein soudin \VUgras{kelluvassa telakassa} \textsuperscript{(53)}, johon
\VUgras{veneet} \textsuperscript{(46)} tulevat hakemaan suojaa, ja minulle oli
ilo käsitellä \VUgras{airoa} \textsuperscript{(47)}. Kiipesin
\VUgras{kannelle} \textsuperscript{(48)} \VUgras{purjeveneessä}
\textsuperscript{(49)}, kiipesin \VUgras{köysiä} \textsuperscript{(52)} pitkin
\VUgras{mastoon} \textsuperscript{(51)} aina \VUgras{raakapuihin}
\textsuperscript{(50)} asti, sitten toistin retkeni \VUgras{jollassa}
\textsuperscript{(54)} raskaiden \VUgras{kalastusveneiden}
\textsuperscript{(55)} välissä, joissa \VUgras{verkot} \textsuperscript{(56)}
kuivuvat.

%%K t10-17-1 p
17. --- \VUgras{Laiturilla} \textsuperscript{(58)} ohitseni kulki mitä
erilaisimpia \VUgras{tavaroita} \textsuperscript{(59)},
\VUgras{laatikoissa} \textsuperscript{(60)} tai \VUgras{paaleissa}
\textsuperscript{(61)}. Ne muodostivat \VUgras{lastin} \textsuperscript{(63)}
\VUgras{purjealuksissa}, ja voimakkaat \VUgras{nosturit}
\textsuperscript{(62)} nostivat ne ylös ja laskivat ne
\VUgras{kuormavaunuihin} \textsuperscript{(64)}, samalla kun \VUgras{ajurit}
ilmoittivat ne ja maksoivat maksut \VUgras{tullikamarissa}
\textsuperscript{(65)}.

%%K t10-18-1 p
18. --- Lopuksi, kun tulin \VUgras{kääntösillalle} \textsuperscript{(66)} tai
\VUgras{sulkuportille} \textsuperscript{(69)}, kulkiessani \VUgras{ruoppaajan}
\textsuperscript{(68)} ohi, joka puhdistaa \VUgras{kanavaa}
\textsuperscript{(67)}, nousin maihin laiturille \VUgras{semaforin}
\textsuperscript{(70)} viereen.

%%K t10-19-1 p
19. --- Tuon laiturin toisella puolella laskevat tietenkin rantaan
\VUgras{sluupit} \textsuperscript{(71)}, jotka tuovat laivaston upseerit maihin.
On ihailtava näytelmä nähdä merimiesten nostavan aironsa tahdissa, samalla kun
vene, jota \VUgras{aalto} \textsuperscript{(72)} kannattaa, kevyesti laskee
maihinnousuportaille. Sillä paikalla voi paremmin tarkkailla \VUgras{vuorovettä}
ja \VUgras{nousuveden} ja \VUgras{laskuveden} liikettä \VUgras{rannalla}
\textsuperscript{(73)}.

%%K t10-tit-7 sub
\VUsaut{3.0mm}
\VUpk{10.2pt}{40}{Seurahuone.}
\VUsaut{3.0mm}

%%K t10-20-1 p
20. --- Kotiin päästäkseni käytin \VUgras{sähköraitiovaunua}
\textsuperscript{(74)}, jonka \VUgras{pysäkki} \textsuperscript{(75)} on sen
\VUgras{pylvään} luona, joka on upseerien seurahuoneen edessä.

%%K t10-20-2 p
Seurahuoneen \VUgras{parvekkeelta} \textsuperscript{(76)}, jota koristavat
\VUgras{ankkurit} \textsuperscript{(77)}, tykit ja kuulat, näin usein loistavan
kokoelman meriupseereja: \VUgras{luutnantteja} \textsuperscript{(78)}
miekkoineen, \VUgras{tykistön kapteeneja} \textsuperscript{(79)} ja
\VUgras{jalkaväen} \textsuperscript{(80)} kapteeneja \VUgras{sapeleineen}.
Heidän takanaan seisoi \VUgras{matruusi} \textsuperscript{(81)}, joka kantoi
heille \VUgras{kaukoputkea}, kun he halusivat erottaa, mitä kaukana tapahtuu.

%%K t10-21-1 p
21. --- Näin myös nuoria \VUgras{aliluutnantteja} \textsuperscript{(82)}, jotka
ottivat vastaan käskyjä \VUgras{komentajaltaan} \textsuperscript{(83)} tai
\VUgras{amiraalilta} \textsuperscript{(84)}, samalla kun \VUgras{pursimies}
\textsuperscript{(85)} odotti viedäkseen heidät alukselle.

%%K t10-22-1 p
22. --- Tuolta parvekkeelta näköala on komea: näkee koko rannan
\VUgras{telttoineen} \textsuperscript{(86)} ja \VUgras{kylpykoppeineen}
\textsuperscript{(87)}, ja näkee, kylpyhetkellä, \VUgras{kylpijät}
\textsuperscript{(88)}, jotka iloisesti leikkivät \VUgras{kasinon}
\textsuperscript{(89)} edessä.

%%K t10-23-1 p
23. --- Rannan päässä rannikko muodostaa pienen kallioisen \VUgras{niemimaan}
\textsuperscript{(91)}, jolle \VUgras{tyrsky} \textsuperscript{(92)} tulee ja
murtuu ja jolle on rakennettu \VUgras{majakka} \textsuperscript{(93)}, joka yöllä
näyttää merenkulkijoille tien. Tuo niemimaa on yhdistetty maahan vain hyvin
kapealla \VUgras{kannaksella} \textsuperscript{(90)}.

%%K t10-tit-8 sub
\VUsaut{3.0mm}
\VUpk{10.2pt}{40}{Yleiskuva rannikoista.}
\VUsaut{3.0mm}

%%K t10-24-1 p
24. --- \VUgras{Jyrkkä rannikko} \textsuperscript{(94)} ulottuu kauas, meren
repimänä, joka oikukkaasti piirtää siihen rivin \VUgras{lahtia}
\textsuperscript{(95)} ja \VUgras{niemiä} (kapeja), mikä tekee siitä hyvin
maalauksellisen.

%%K t10-24-2 p
\VUgras{Ylätasangolla} \textsuperscript{(96)}, joka hallitsee
\VUgras{(meri)salmea} \textsuperscript{(98)}, on hyvin tärkeä \VUgras{linnake}
\textsuperscript{(97)} ja muutamia « huviloita ».

%%K t10-25-1 p
25. --- Tuo salmi erottaa mantereesta hyvin vaarallisen \VUgras{saaren}
\textsuperscript{(99)}, jota ympäröivät \VUgras{karit} \textsuperscript{(100)}
ja \VUgras{tyrskyt}, joilla veneet ja jopa suuret alukset joskus ajavat
karille. Avun nopeudesta ja \VUgras{pelastusveneiden} nopeasta saapumisesta
huolimatta \VUgras{haaksirikot} ovat hirvittäviä, ja päiväntasauksen
\VUgras{myrskyjen} aikana useita aluksia on hukkunut väkineen ja tavaroineen.

%%K t10-25-2 p
Tuosta naapuruudesta huolimatta merenkulkijat käyvät \VUgras{satamassa}
paljon.
\VUsaut{6.0mm}
\VUfilet{22mm}
